% Part: incompleteness
% Chapter: representability-in-q
% Section: c

\documentclass[../../../include/open-logic-section]{subfiles}

\begin{document}

\olfileid{inc}{req}{cee}
\olsection{تابع‌های $C$}

بگذارید $C$ کوچک‌ترین مجموعهٔ تابع‌هایی باشد که موارد زیر را در بر دارد:
\begin{enumerate}
\item صفر؛
\item جانشین؛
\item جمع؛
\item ضرب؛
\item تصویرها؛ و
\item تابع مشخصهٔ برابری، $\Char{=}$؛
\end{enumerate}
و تحت عملیات زیر بسته است:
\begin{enumerate}
\item ترکیب؛ و
\item جست‌وجوی نامحدود، وقتی بر تابع‌های منظم اعمال شود.
\end{enumerate}
به یاد آورید که محدودیت آخر فقط به این معناست که عمل $\mu$ را تنها
وقتی می‌توان به کار برد که نتیجه تابعی تام باشد. این را با تعریف
تابع‌های \emph{بازگشتی عمومی} مقایسه کنید: اینجا جمع، ضرب و
$\Char{=}$ را افزوده‌ایم، اما بازگشت اولیه را کنار گذاشته‌ایم.

آشکار است که هر تابع در $C$ بازگشتی است، زیرا جمع، ضرب و $\Char{=}$
بازگشتی‌اند. نشان خواهیم داد عکس این گزاره نیز صادق است؛ یعنی با دیگر
عملیات موجود در $C$ می‌توانیم بازگشت اولیه را انجام دهیم.

\end{document}
